\chapter{MLOps und der Modell-Lebenszyklus}

\label{chap:mlops_modelllebenszyklus}

% ============================================================
% LERNZIELE
% ============================================================
\begin{lernziele}
Nach diesem Kapitel kannst du:
\begin{itemize}
    \item den Unterschied zwischen klassischem DevOps und LLMOps erklären
    \item ein Versionierungs- und Release-Setup für Prompts, Modelle und Datasets entwerfen
    \item eine CI/CD-Pipeline für LLM-Änderungen konfigurieren
    \item Produktionsdrift und Kostenabweichungen systematisch erkennen
    \item Verantwortlichkeiten, Governance und Audit für KI-Produkte strukturieren
\end{itemize}
\end{lernziele}

% ============================================================
% MOTIVATION
% ============================================================
\section{Motivation}

LLM-Anwendungen bestehen aus mehr als Code. Sie kombinieren Modelle, Daten, Prompts und Infrastruktur. Deshalb reicht klassisches Software-Engineering nicht aus: Es muss ein Lebenszyklus entstehen, der alle Komponenten versioniert, testet, deployed und überwacht.

MLOps ist die Disziplin, die diese Kette operationalisiert. Für AI Engineers ist sie entscheidend, weil sie aus einem einmaligen Prototypen eine stabile, skalierbare Produktionsanwendung macht.

% ============================================================
% GRUNDLAGEN
% ============================================================
\section{Was macht MLOps anders?}

Im Vergleich zur klassischen Softwareentwicklung sind drei Aspekte zentral:

\begin{enumerate}[leftmargin=*]
    \item \textbf{Nicht-deterministisches Verhalten:} Der gleiche Prompt liefert nicht immer denselben Output. Deshalb müssen Tests probabilistisch und robust gestaltet werden.
    \item \textbf{Externe Modellabhängigkeit:} Viele Anwendungen nutzen Managed APIs. Verfügbarkeit, Latenz und Kosten liegen oft außerhalb der eigenen Kontrolle.
    \item \textbf{Daten- und Modelldrift:} Das reale Nutzerverhalten ändert sich, und der promptbasierte Code altert schneller als klassischer Application Code.
\end{enumerate}

% ============================================================
% ARCHITEKTUR DES LEBENSZYKLUS
% ============================================================
\section{Der LLMOps-Lebenszyklus}

Eine LLMOps-Architektur besteht aus fünf Hauptphasen:

\begin{description}[leftmargin=*,style=nextline]
    \item[Entwickeln] Prompts, Modell-Konfigurationen, Daten-Pipelines und Service-Code erstellen
    \item[Testen] Golden Datasets, Prompt-Tests, Kosten- und Latenzprüfungen ausführen
    \item[Ausrollen] Versionierte Releases, Canary-Deployments und Rollbacks verwalten
    \item[Überwachen] Drift, Feedback, Kosten und Performance in Produktion beobachten
    \item[Verbessern] Regeln, Daten und Prompts iterativ anpassen und den Kreislauf erneuern
\end{description}

\subsection{Systemkomponenten}

\begin{itemize}[leftmargin=*]
    \item \textbf{Versionierung}: Git für Code und Prompts, DVC/LakeFS für Trainingsdaten, Model Cards für Modelle
    \item \textbf{CI/CD}: Automatisierte Tests, Review-Gates, Canary-Deployments
    \item \textbf{Observability}: Metriken, Tracing, Log-Events, Drift-Metriken
    \item \textbf{Governance}: Verantwortlichkeiten, Dokumentation, Audit-Logs
    \item \textbf{Feedback}: Nutzerfeedback, menschliche Evaluierung, Incident Reports
\end{itemize}

% ============================================================
% VERSIONIERUNG
% ============================================================
\section{Versionierung von Prompts, Daten und Modellen}

\subsection{Prompts als Code}
Prompts gehören in die Codebasis. Jede Änderung muss eine Git-Historie haben, eine Review erhalten und bei Bedarf zurücksetzbar sein. Dazu gehören:

\begin{itemize}[leftmargin=*]
    \item \textbf{Prompt-Templates} als parametrische Dateien
    \item \textbf{Prompt-Tests} gegen Golden Datasets
    \item \textbf{Release Notes} für Prompt-Änderungen
\end{itemize}

\subsection{Datenversionierung}

Daten für Evaluationen, Referenzen und semantische Indizes müssen versioniert werden. Für AI Engineers gelten folgende Regeln:

\begin{itemize}[leftmargin=*]
    \item \textbf{Daten-Snapshots} für historische Vergleichbarkeit
    \item \textbf{Schema-Checks} zur Erkennung von Drift
    \item \textbf{Reproduzierbare Builds} durch festgehaltene Datenversionen
\end{itemize}

\subsection{Modelldokumentation}

Jedes Modell braucht eine \textbf{Model Card}. Diese beschreibt:

\begin{itemize}[leftmargin=*]
    \item Zweck und Einsatzgebiet
    \item Trainingsdaten- oder Prompt-Basis
    \item bekannte Schwächen und Risiken
    \item Performance-Metriken und Kosten
\end{itemize}

% ============================================================
% CICD
% ============================================================
\section{CI/CD für LLM-Systeme}

Eine LLMOps-Pipeline enthält neben standardmäßigen Code-Checks auch prompt- und datenorientierte Schritte.

\subsection{Beispiel-Pipeline}

\begin{itemize}[leftmargin=*]
    \item \textbf{Linting}\ für Code, Prompts und JSON-Schema
    \item \textbf{Golden-Dataset-Tests}\ zur Qualitätsabsicherung
    \item \textbf{LLM-Judge-Tests}\ zur semantischen Bewertung
    \item \textbf{Kosten-/Latenz-Checks}\ zur Budgetkontrolle
    \item \textbf{Canary Deploy}\ zur risikoreduzierten Einführung
\end{itemize}

\subsection{Praktischer Ablauf}

Eine CI/CD-Stage kann so aussehen:

\begin{lstlisting}[language=Bash,caption={Beispiel-Workflow für LLMOps in GitHub Actions}]
name: LLMOps Pipeline
on: [pull_request, push]

jobs:
  test:
    runs-on: ubuntu-latest
    steps:
      - uses: actions/checkout@v4
      - name: Install dependencies
        run: pip install -r requirements.txt
      - name: Prompt lint
        run: python scripts/prompt_lint.py prompts/
      - name: Run golden dataset tests
        run: python scripts/evaluate_golden.py
      - name: Validate JSON schema
        run: python scripts/validate_output_schema.py

  canary:
    needs: test
    if: github.event_name == 'push' && github.ref == 'refs/heads/main'
    runs-on: ubuntu-latest
    steps:
      - uses: actions/checkout@v4
      - name: Deploy canary
        run: ./scripts/deploy_canary.sh
      - name: Run canary smoke tests
        run: python scripts/canary_smoke_test.py
\end{lstlisting}

% ============================================================
% MONITORING
% ============================================================
\section{Monitoring und Drift-Erkennung}

Überwachung bedeutet nicht nur, dass ein Dienst erreichbar ist. Für LLM-Systeme zählen auch Inhalt und Qualität.

\subsection{Wichtige Produktionsmetriken}

\begin{description}[leftmargin=*,style=nextline]
    \item[Qualitäts-Drift] Veränderung der Antwortqualität über Zeit. Gemessen durch Stichproben, Feedback oder automatisierte Judges.
    \item[Daten-Drift] Veränderungen in den Eingabedaten (z. B. Länge, Domäne, Token-Verteilung).
    \item[Kosten-Drift] Unerwartete Steigerungen durch höhere Tokens, längere Ausgaben oder neue Nutzerprofile.
    \item[Performance-Drift] Steigende Latenzen, höhere Fehlerquoten oder vermehrte Timeouts.
\end{description}

\subsection{Alerting}

Alerts sollten nicht nur bei Ausfällen ausgelöst werden, sondern auch bei:

\begin{itemize}[leftmargin=*]
    \item \textbf{Signifikanter Kostenabweichung} im Vergleich zur Baseline
    \item \textbf{Spürbarer Qualitätsverschlechterung} im Golden Dataset
    \item \textbf{Datenanomalien} bei Eingabeverteilungen
    \item \textbf{Anstieg von User Complaints} oder Ablehnungen
\end{itemize}

% ============================================================
% GOVERNANCE
% ============================================================
\section{Governance und Verantwortlichkeiten}

Eine stabile KI-Produktion braucht klare Rollen:

\begin{itemize}[leftmargin=*]
    \item \textbf{Model Owner}: Verantwortlich für die Modellwahl, Release-Planung und Dokumentation
    \item \textbf{Data Steward}: Sorgt für Datenqualität, Versionierung und Compliance
    \item \textbf{Product Owner}: Bewertet Nutzerfeedback und priorisiert Verbesserungen
    \item \textbf{Security Owner}: Überprüft Zugriff, Datenschutz und Missbrauchsschutz
\end{itemize}

\subsection{Audit und Dokumentation}

Jede Modelländerung sollte nachvollziehbar sein. Wichtige Artefakte sind:

\begin{itemize}[leftmargin=*]
    \item Release Notes für Prompts, Modelle und Indexdaten
    \item Logs von Evaluationen und Governance Reviews
    \item Model Cards und Risikoanalysen
    \item Approval-Gates vor Deployments
\end{itemize}

% ============================================================
% ZUSAMMENFASSUNG
% ============================================================
\begin{zusammenfassung}
MLOps macht LLM-Anwendungen produktionsreif: nicht nur durch Deployment, sondern durch Versionierung, Tests, Monitoring und Governance. Ein strukturierter Lebenszyklus reduziert Risiko, Kosten und Drift, und ermöglicht wiederholbare Verbesserungen.
\end{zusammenfassung}

\begin{uebungen}
    \item Beschreibe einen konkreten Ablauf, wie du ein Prompt-Update von der Review bis zum Canary-Release bringen würdest.
    \item Nenne drei Drift-Indikatoren, die für eine Conversational AI besonders wichtig sind.
    \item Erkläre, warum eine Model Card für API-basierte Modelle trotzdem sinnvoll ist.
\end{uebungen}
